\section{一个计算，多种形态：有界前缀截断点积}
\label{sec:case-study}

在研究编译器怎样改变程序之前，必须先说清程序\emph{做什么}。本章不讨论寄存器、基本块或
目标文件，而把案例压缩成一个数学函数和几条不变量。它们是后续所有表示共同承担的语义线索。

\subsection{规范源程序}

代码清单~\ref{lst:sysy-case} 是全文唯一的源级案例。SysY 运行时预先提供 \texttt{getint}、
\texttt{putint} 与 \texttt{putch}；数组形参传递首元素地址，标量形参按值传递。这一区别看似
只是语法规则，后来会分别成为 LLVM 的指针/整数参数以及 RISC-V ABI 中的寄存器值。

\begin{lstlisting}[language=C,caption={规范性 SysY 案例。},label={lst:sysy-case}]
int clipped_dot(int a[], int b[], int n, int lo, int hi) {
  int acc = 0;
  int i = 0;
  while (i < n) {
    int term = a[i] * b[i];
    if (term < lo) {
      term = lo;
    } else {
      if (term > hi) {
        term = hi;
      }
    }
    acc = acc + term;
    i = i + 1;
  }
  return acc;
}

int main() {
  int a[8] = {-4, 1, 7, 3, 9, -2, 6, 5};
  int b[8] = { 2,-3, 1, 4,-1,  5, 2, 3};
  int n = getint();
  if (n < 0) {
    n = 0;
  } else {
    if (n > 8) {
      n = 8;
    }
  }
  putint(clipped_dot(a, b, n, -8, 12));
  putch(10);
  return 0;
}
\end{lstlisting}

令
\[
  \operatorname{clip}(x,l,h)=\min(h,\max(l,x)),\qquad
  k=\operatorname{clip}(n,0,8).
\]
程序实现的函数为
\begin{equation}
  D(n)=\sum_{i=0}^{k-1}\operatorname{clip}(a_i b_i,-8,12),
  \label{eq:clipped-dot}
\end{equation}
其中空和等于 0。两个固定数组的乘积
\[
(-8,-3,7,12,-9,-10,12,15)
\]
经逐项截断变为
\[
(-8,-3,7,12,-8,-8,12,12),
\]
所以长度 $0$ 到 $8$ 的前缀和依次是 $0,-8,-11,-4,8,0,-8,4,16$。这些数字的主要作用
不是构造一个庞大的测试仪式，而是让读者随时能够把抽象机制重新落到一个可心算的计算上。

% Worked mathematical case / 数学案例总览
\begin{figure}[htbp]
  \centering
  \begin{tikzpicture}
    \begin{axis}[
      name=clipplot,at={(0,0)},anchor=south west,
      width=4.30cm,height=4.25cm,
      xmin=-15,xmax=18,ymin=-12,ymax=15,
      axis lines=middle,ticks=none,clip=false,
      xlabel={$x$},xlabel style={font=\footnotesize},
      every axis plot/.append style={semblue,very thick}]
      \addplot coordinates {(-15,-8) (-8,-8) (12,12) (18,12)};
      \node[vis/label,anchor=north] at (axis cs:-11,-8) {$\ell=-8$};
      \node[vis/label,anchor=south,rotate=38] at (axis cs:2,2) {恒等区间};
      \node[vis/label,anchor=north] at (axis cs:15,12) {$h=12$};
    \end{axis}
    \node[vis/label,anchor=south west,font=\bfseries\footnotesize]
      at (clipplot.north west) {(a) $y=\operatorname{clip}(x,-8,12)$};

    \begin{scope}[shift={(4.82cm,.05cm)}]
      \node[vis/label,anchor=west,font=\bfseries\footnotesize] at (0,3.72) {(b) 八项计算与前缀和};
      \node[vis/label,anchor=east] at (0,3.05) {$a_i b_i$};
      \node[vis/label,anchor=east] at (0,2.35) {截断后};
      \foreach \x/\raw/\val/\chg in {
        .45/-8/-8/0,1.28/-3/-3/0,2.11/7/7/0,2.94/12/12/0,
        3.77/-9/-8/1,4.60/-10/-8/1,5.43/12/12/0,6.26/15/12/1}
      {
        \node[vis/artifact,minimum width=.72cm,minimum height=.48cm] at (\x,3.05) {$\raw$};
        \ifnum\chg=1
          \node[vis/decision,minimum width=.72cm,minimum height=.48cm] at (\x,2.35) {$\val$};
        \else
          \node[vis/artifact,minimum width=.72cm,minimum height=.48cm] at (\x,2.35) {$\val$};
        \fi
      }
      \draw[->,black!55] (.12,.20) -- (6.7,.20) node[right,vis/label] {$k$};
      \draw[->,black!55] (.12,.20) -- (.12,1.85) node[above,vis/label] {$D(k)$};
      \draw[vis/semantic,line width=1.1pt]
        plot coordinates {(.12,1.15) (.45,.72) (1.28,.56) (2.11,.93)
          (2.94,1.58) (3.77,1.15) (4.60,.72) (5.43,1.36) (6.26,1.80)};
      \foreach \x/\y in {.12/1.15,.45/.72,1.28/.56,2.11/.93,
          2.94/1.58,3.77/1.15,4.60/.72,5.43/1.36,6.26/1.80}
        \filldraw[fill=white,draw=semblue,line width=.8pt] (\x,\y) circle (1.15pt);
      \foreach \x/\k in {.12/0,.45/1,1.28/2,2.11/3,2.94/4,
          3.77/5,4.60/6,5.43/7,6.26/8}
        \node[vis/label,below] at (\x,.20) {$\k$};
      \node[vis/note,anchor=west] at (.18,-.47)
        {$k=\operatorname{clip}(n,0,8)$；橙色格表示截断改变了乘积。};
    \end{scope}
  \end{tikzpicture}
  \caption{案例计算的两次限制发生在不同对象上：输入决定前缀长度 $k$，逐项截断决定每个加数；蓝色折线给出 $k=0,\ldots,8$ 时的前缀和。}
  \label{fig:worked-case}
\end{figure}


图中的同一乘法在各层换了身份：源程序里它是两个下标表达式；带类型 AST 中还包含隐式的
左值到值转换；LLVM IR 把元素地址、加载与乘法拆开；RV64 用地址算术、\texttt{lw} 和
\texttt{mulw} 实现；目标文件只保存编码后的字节和必要的符号关系。跨层延续的不是变量名或
排版，而是“取两个第 $i$ 个整数并相乘”这条数据依赖。

\subsection{三条不变量连接函数与循环}

案例的第一个不变量是内存安全：
\[
  0\le i<k\le 8.
\]
输入先被限制到 $[0,8]$，循环条件成立后才会求值数组下标，因此每次访问都落在两个数组内部。
这让后文可以专注于地址是怎样表示和优化的，而不必把偶然越界行为混入讨论。

第二个不变量把命令式循环连接到式~\eqref{eq:clipped-dot}。记
\(c_j=\operatorname{clip}(a_jb_j,-8,12)\)。在闭合的 \texttt{main} 调用中，传给
\texttt{clipped\_dot} 的长度正是 $k$；循环头应满足
\begin{equation}
  I(i,acc)\;\equiv\;
  0\le i\le k\;\land\;
  acc=\sum_{j=0}^{i-1}c_j.
  \label{eq:loop-invariant}
\end{equation}
初始化时 $i=0,acc=0$，空和使 $I(0,0)$ 成立；若 $I(i,acc)$ 成立且 $i<k$，一次循环执行
\((i,acc)\mapsto(i+1,acc+c_i)\)，于是新状态仍满足 $I$；退出时条件否定，结合
$0\le i\le k$ 得到 $i=k$，因而 $acc=D(n)$。这三个步骤不是为了把文章改写成证明报告，
而是给后文一根真正守恒的轴：SysY 的两个赋值、SSA 循环头的两个 \texttt{phi}、RV64 的回边
以及 MLIR 的归约迭代参数，都在编码同一个状态转移。

第三个不变量是整数范围。原始乘积落在 $[-10,15]$，截断后的每项落在 $[-8,12]$，八项和必在
$[-64,96]$。因此本例的 32 位有符号乘法和加法不会溢出。这个范围证明虽短，却有直接的工程
后果：优化器无需借助有符号溢出的未定义行为来解释本例，读者也可以把源级整数与 LLVM
\texttt{i32} 运算对应起来。若数组改为任意用户输入，同一实现就需要更宽的中间类型、显式
溢出策略或更严格的前置条件；追求更一般的接口，会把额外语义成本传到每一层。

逐项截断还揭示了一个容易被“点积”名称遮蔽的区别。这里先计算每个
$\operatorname{clip}(a_i b_i,-8,12)$，再用普通加法归约；由于和不会溢出，改变加法树的形状
不会改变结果。若改为每次更新累加器后都做饱和截断，混合正负项会破坏结合性，树形归约和
向量化便不能直接保持顺序语义。编译器从不优化“看起来像点积”的东西；它优化满足特定代数
条件的运算。

\subsection{为何一个小程序足以打开整个系统}

案例的每个结构都向后投下一道影子：
\begin{itemize}
  \item \texttt{a[i]} 与 \texttt{b[i]} 要经过类型检查，随后变成元素地址、加载以及目标机寻址；
  \item \texttt{i} 与 \texttt{acc} 在未优化表示中可能位于栈槽，在 SSA 中则成为循环携带值，
        最终竞争有限的物理寄存器；
  \item 两级截断条件先成为控制流图中的分支与汇合，也可能被优化成选择或最值序列；
  \item 五参数函数迫使后端服从调用约定，即使另一种局部寄存器安排看上去更漂亮；
  \item \texttt{putint} 从带类型调用变为外部符号和重定位，直到链接时才获得定义。
\end{itemize}
这个选择体现了一条写作原则：案例无需模拟真实应用的规模，只需让重要机制不可回避。加入递归、
浮点或并发会产生新的叙事线，却不会使当前这条线解释得更清楚。

\subsection{四种材料，而非四个彼此竞争的程序}

后文会看到四种材料。规范 SysY 源文件给出源级含义；C17 观察载体仅为借助成熟 Clang 展示
预处理、记号与 AST，在本案例的整数子集内镜像同一算法；Clang 生成的 IR 和汇编展示生产编译器
作出的选择；手写 LLVM IR 与 RV64 汇编则让我们从中间层直接阅读同一计算。它们服务于同一
讲解，却承担不同角色，不能把“Clang 接受这份 C”误写成“Clang 是一个完整 SysY 前端”。

课程实现还把这四种形式在 $n=-1,0,1,4,8,10$ 上运行，所得输出依次为
\texttt{0,0,-8,8,16,16}。这一小表足以确认配套产物可以实际汇编、链接和运行；它不是文章的
中心，更不替代对表示与机制的解释。确切命令、工具版本和运行时诊断保存在
\texttt{preflight/README.md} 中。

\subsection{路线图：每一层得到什么，又放弃什么}

从源程序向机器前进，并非给每个变量不断换一种拼写。前端保留声明、类型和源位置，适合诊断，
却不适合全局数据流优化；LLVM IR 放弃大部分源语法，以基本块、显式内存操作和 SSA 关系换取
跨语言优化能力；机器表示接受寄存器类别、指令延迟和 ABI 约束，换取可执行性；ELF 又把已确定
的字节与尚未确定的地址方程一起封装，换取独立编译和部署弹性。

因此，本教程式旅程的核心并不是“源代码经过了多少步骤”，而是三个连续问题：当前表示让哪种
推理变得容易？它不再保存什么？被推迟的决定将由谁完成？下一章从第一个边界开始：字符与宏
本身并不知道 \texttt{a[i]} 合不合法，前端必须把线性记号构造成带绑定、类型和值类别的程序。
